% BEGIN REV09_PELL_MOMENTO_PREVIO_IV
\subsection{Momento orientado de profundidad}
\label{corpus9:iv:pell-momento}

La realización del registro de profundidad con cuarto de giro
define el momento paramétrico
\[
 \mathcal C_2(r)=
   \sum_{k\geq0}\frac{(-1)^kr^{2k+1}}{(2k+1)^2},
 \qquad 0<r<1,\qquad
 \mathcal D=r\frac{d}{dr}.
\]
La convergencia local uniforme permite derivar término a término:
\[
 \mathcal D\mathcal C_2(r)=\arctan r,\qquad
 \mathcal D^2\mathcal C_2(r)=\frac{r}{1+r^2}.
\]
En efecto, la primera serie derivada es la integral desde cero
de la serie geométrica \(1/(1+r^2)\). Una segunda derivación
logarítmica produce la última igualdad. Las condiciones
\(\mathcal C_2(0)=\mathcal D\mathcal C_2(0)=0\) fijan las
constantes de integración y dan
\[
 \mathcal C_2(r)=\int_0^r\frac{\arctan t}{t}\,dt
   =\int_0^r\frac{\log(r/t)}{1+t^2}\,dt.
\]
La dominación por \(\sum_{n\geq1}n^{-2}\) conserva el límite
\(\mathcal C_2(1)=G_{\mathrm{Cat}}\). La segunda derivada
en \(r=1\) se entiende mediante su límite de Abel.
% END REV09_PELL_MOMENTO_PREVIO_IV
